\chapter{轨迹族、界面通量与机理图谱}
\label{ch:atlas}

交叉表给出中位。要把中位读回过程，需要同一初值上的轨迹族：四个策略、三种厚度、三种档位，以及界面通量、进度分量和路径代理随时间的形状。下面的轨迹全部来自种子 7、肉糜、初温 \(8^\circ\mathrm{C}\)、含水 0.72。温度场图由十六条剖面按 \((d/L)^{1.15}\) 插值着色，横向没有第二维梯度。演示页面默认读这一条，附录\ref{app:trace} 给出同一条的抽样表。

\section{四条策略}

图\ref{fig:trpol}给出开环、腔温比例、无探针与对偶的冷点、进度和加热。开环在 1170 s 切断，冷点 \(16^\circ\mathrm{C}\) 量级，进度 0.258，偏差 \(-0.849\)，方差留在先验 0.64。腔温比例把加热开到 3300 s，能量 3312，进度停在上沿 \(+0.150\)，方差仍是 0.64。无探针没有蒸汽窗，能量 2319，进度同样走到上沿，方差仍是 0.64。对偶在 192 s 把方差降到 0.224，1866 s 停机，偏差 \(+0.145\)，能量 1864。

\begin{figure}[htbp]
\centering
\includegraphics[width=0.90\textwidth]{X_trace_policy.png}
\caption{种子 7、30 mm、组合档：四条策略的冷点、进度与加热。只有对偶把方差降到 0.224。}
\label{fig:trpol}
\end{figure}

图\ref{fig:d1t}单独画出表面、冷点与估计。估计从负值爬起，约 5 min 越过冷点，在停机前仍高于冷点，切入保持后落到冷点之下。跟踪器给出带滞后的点估计：第一个记录点出现 \(-7.92^\circ\mathrm{C}\)，来自先验热阻 5 的滞后上限；1200 s 时估计 \(61.27^\circ\mathrm{C}\) 高于冷点 \(46.65^\circ\mathrm{C}\)。图\ref{fig:d6v}给出方差：探针窗结束时从 0.64 落到 0.224，之后不再变化。冷点进度在 25--30 min 饱和到 1，褐变代理几乎不离开零——家用组合档把表面湿度维持得较高，褐变通道几乎不动作。

\begin{figure}[htbp]
\centering
\includegraphics[width=0.82\textwidth]{D1_temperature.png}
\caption{种子 7、对偶、30 mm：表面、冷点与估计。估计可以高于冷点。}
\label{fig:d1t}
\end{figure}

\begin{figure}[htbp]
\centering
\includegraphics[width=0.82\textwidth]{D6_variance.png}
\caption{种子 7、对偶：热阻方差。窗后从 0.64 落到 0.224，之后不变。}
\label{fig:d6v}
\end{figure}

无探针在这一条 30 mm 轨迹上走到了上沿，S1 无探针的中位却是 \(-0.467\)。演示轨迹取初温 \(8^\circ\mathrm{C}\)，S1 含 \(5^\circ\mathrm{C}\) 的一半样本，并且混有 15 mm 与 50 mm。单条走到上沿，全表中位仍由厚度混合决定。对偶在这一条上也走到上沿，S1 对偶中位 \(-0.313\)，同样是厚度混合的结果。把演示页默认条件当成全表中位，会把 30 mm、\(8^\circ\mathrm{C}\) 的一次成功读成“对偶已经越过设计线”。图\ref{fig:f12}把档位研究里若干轨迹叠在一起，形状与种子 7 一致：厚件进度晚离开零，薄件早饱和，探针窗在近表面留下短暂高温。叠图里没有一条 50 mm 曲线在 55 min 内贴近 0.85。

\begin{figure}[htbp]
\centering
\includegraphics[width=0.86\textwidth]{F12_traces.png}
\caption{档位研究中的冷点与进度轨迹族。厚件进度晚离开零。}
\label{fig:f12}
\end{figure}

\section{三种厚度}

图\ref{fig:trthick}固定对偶、组合档，打开 15、30、50 mm。15 mm 在时域内把进度推到 1，停止 3300 s，偏差 \(+0.150\)，能量 1621。30 mm 即上节的对偶轨迹。50 mm 在 3300 s 进度只有 0.068，偏差 \(-0.782\)，能量 3310。三条结束方差都是 0.224：探针在三种厚度上都接受了新热阻。厚件的负偏差在方差已经下降之后仍然在。

\begin{figure}[htbp]
\centering
\includegraphics[width=0.90\textwidth]{X_trace_thick.png}
\caption{种子 7、对偶、组合档：15、30、50 mm 的进度。50 mm 在 55 min 内停在 0.068。}
\label{fig:trthick}
\end{figure}

15 mm 的停止时刻是时域终点，因为进度很早就饱和，停机判据在薄件上被上沿和最小时间 240 s 之外的保持逻辑共同改写。日志里 15 mm 对偶的加热在中后段已经降到低功率，进度贴着 1。\(e_{\alpha}=+0.150\) 是 \(\alpha_m=1\) 相对 0.85 的差，饱和之后进度轴没有分辨率。

\section{三种档位}

图\ref{fig:trcook}固定对偶、30 mm，打开蒸、烤、组合。蒸档加热上限 0.25，能量 828，进度 0.021，偏差 \(-0.829\)，方差留在 0.64：加热被夹住之后，探针窗的表面温差不够，热阻更新没有接受。烤档关闭蒸汽，能量 2353，进度走到上沿，方差 0.224。组合档即演示轨迹。蒸档失败在允许集合，失败在探针没有买到方差下降，两件事同时发生。

\begin{figure}[htbp]
\centering
\includegraphics[width=0.90\textwidth]{X_trace_cook.png}
\caption{种子 7、对偶、30 mm：蒸、烤、组合。蒸档进度停在 0.021，方差留在 0.64。}
\label{fig:trcook}
\end{figure}

档位研究里对偶蒸/烤/组合的中位是 \(-0.830\)/\(+0.150\)/\(+0.146\)，与这三条轨迹的符号一致。烤档与组合档的中位被薄件抬到正值；50 mm 在烤档上同样走不到设计线。

\section{厚度--时间场与插值色带}

图\ref{fig:xtx}把种子 7 对偶组合档的十六条剖面沿时间展开。探针窗在近表面留下短暂高温条带，冷点一侧在 10 min 之内几乎不动。28 min 附近冷点越过 \(70^\circ\mathrm{C}\)，此时仍是品质加热，请求 0.476。演示页面的着色器读的是这一帧的表面温度与冷点两个标量，中间格子按 \(u^{1.15}\) 插值，\(u\) 是归一化深度。横向没有第二维梯度。指数 1.15 是显示约定，不是第二套物理。

\begin{figure}[htbp]
\centering
\includegraphics[width=0.88\textwidth]{X_field_xt.png}
\caption{种子 7 对偶组合档的厚度--时间温度场。探针窗的高温条带停在近表面。}
\label{fig:xtx}
\end{figure}

含水场见图\ref{fig:moist}。平均含水从 0.72 降到约 0.65，表面格子降得更快。进度 \(\alpha_m\) 在 20 min 之后才离开零，表面褐变代理 \(\alpha_b\) 随表面温度更早升起。二者对同一蒸汽指令的符号不必相同：探针窗增加近表面湿度、压低褐变，同时给冷点的潜热增量几乎为零。

\begin{figure}[htbp]
\centering
\includegraphics[width=0.82\textwidth]{X_moist.png}
\caption{种子 7：平均含水与表面含水。结壳门槛 0.45 在 55 min 内未触及。}
\label{fig:moist}
\end{figure}

种子 7 平均含水从 0.72 降到约 0.65，结壳门槛 0.45 在 55 min 内未触及。质量通道开或关不改写跟踪器，第\ref{ch:ablate}章把这一通道写成零效应；含水场本身仍随蒸发缓慢下降，下降量进入界面活度的饱和段，停机偏差对三档初含水仍然不敏感。

\section{路径代理随时间}

日志里的路径 A、B、C 是该帧通量代理，\eqref{eq:path} 的偏效应积分留在规格里。图\ref{fig:pathst}给出三条代理。探针窗内蒸汽为 1，路径 A 升到 180，路径 B 落到 0：冷凝方向关掉蒸发冷却。窗后蒸汽回到 0.15，路径 A 回到 27，路径 B 回到 1958。路径 C 在 1680 s 之前是 0，之后随结壳变量爬到 \(10^{-2}\) 量级。结壳没有乘回导热，路径 C 的数值几乎不进入温度。图\ref{fig:f15}给出执行器时间序列：蒸汽峰值只出现在探针窗。图\ref{fig:f14cav}给出种子 7 的腔温与腔湿：探针窗内腔湿先降后升，品质段腔温停在 \(160^\circ\mathrm{C}\) 附近，保持段随请求下降。加热请求与指令两条线重合，屏障不动作，因为加热器停在 \(196.5^\circ\mathrm{C}\)，低于 \(230^\circ\mathrm{C}\) 阈值。

\begin{figure}[htbp]
\centering
\includegraphics[width=0.84\textwidth]{F15_actuators.png}
\caption{档位研究执行器时间序列。蒸汽峰值只出现在探针窗。}
\label{fig:f15}
\end{figure}

\begin{figure}[htbp]
\centering
\includegraphics[width=0.82\textwidth]{F14_cavity.png}
\caption{种子 7 腔温与腔湿。品质段腔温停在 \(160^\circ\mathrm{C}\) 附近。}
\label{fig:f14cav}
\end{figure}

\begin{figure}[htbp]
\centering
\includegraphics[width=0.86\textwidth]{X_paths_time.png}
\caption{种子 7 路径代理。探针窗内 A 升、B 落；C 在积分里接近零。}
\label{fig:pathst}
\end{figure}

界面上的三根条若用蒸汽指令和方差拼出，那是显示公式。要把路径写成解释句，需要在后验抽样上替换蒸汽轨迹再积分，符号不稳定时不生成句子。本次没有这一抽样。

\section{表面温差、Antoine 项与算子残差}

图\ref{fig:gap}给出表面与冷点的温差。探针窗内温差升到 \(90\,\mathrm{K}\) 量级，正是 \(r_{\mathrm{proxy}}\) 使用的真值 \(\Delta T\)。实施到机器上必须改成可安装通道。Antoine 项打开与关闭的成对积分已经在第\ref{ch:result}章给出符号随厚度翻转；图\ref{fig:antd}把 216 对的差按厚度画成分布。15 mm 全部为正，50 mm 全部为负。

\begin{figure}[htbp]
\centering
\includegraphics[width=0.82\textwidth]{X_gap.png}
\caption{种子 7 表面与冷点温差。探针窗内温差约 \(90\,\mathrm{K}\)，\(r_{\mathrm{proxy}}\) 读取这一真值。}
\label{fig:gap}
\end{figure}

\begin{figure}[htbp]
\centering
\includegraphics[width=0.82\textwidth]{X_ant_delta.png}
\caption{打开 Antoine 项之后的进度偏差变化。15 mm 上移，50 mm 下移。}
\label{fig:antd}
\end{figure}

两节点教师上训练的主干加分支网络，留出均方误差 1683，对应均方根约 \(41^\circ\mathrm{C}\)。训练集 3276、留出 820、80 个周期。图\ref{fig:oper}给出留出误差。该网络没有接到控制器。均方根与过程温升同量级，外推到十六条场需要另一套教师。

\begin{figure}[htbp]
\centering
\includegraphics[width=0.78\textwidth]{X_operator.png}
\caption{两节点教师网络的留出误差。均方根约 \(41^\circ\mathrm{C}\)，未接入控制器。}
\label{fig:oper}
\end{figure}

有限时域搜索在某一个状态下对 60 s、加热点 0.7 给出预测冷点 \(74.73^\circ\mathrm{C}\)。图\ref{fig:mpc}只画这一个记录点。无探针策略在核里是饱和线性请求\eqref{eq:ureq}。

\begin{figure}[htbp]
\centering
\includegraphics[width=0.72\textwidth]{X_mpc_horizon.png}
\caption{两节点上一次 60 s 搜索的记录点。预测冷点 \(74.73^\circ\mathrm{C}\)。}
\label{fig:mpc}
\end{figure}

同一初值上，只有对偶把方差降到 0.224。蒸档让探针失败：加热上限 0.25 之后窗内表面温差不够，热阻更新不接受。50 mm 在方差已经下降之后仍停在 0.068。路径 C 在积分里接近零，因为结壳门槛 0.45 未触及。Antoine 项的符号随厚度翻转。教师网络的均方根是 \(41^\circ\mathrm{C}\)，该网络没有接到控制器。演示页面默认读组合档、对偶、30 mm 这一条。把默认条件当成 S1 中位，会把 30 mm、\(8^\circ\mathrm{C}\) 的一次成功读成“对偶已经越过设计线”。
